\section{Trabajo mecánico}

\subsection{Definición}
El trabajo realizado por una fuerza constante es el producto de la componente de la fuerza en la dirección del desplazamiento por la magnitud del desplazamiento.
\begin{equation}
    W = \vec{F} \cdot \Delta \vec{x} = F \cdot d \cdot \cos \theta
\end{equation}

\subsection{Unidades}
\begin{itemize}
    \item Joule (J): \(1 \, \text{J} = 1 \, \text{N} \cdot \text{m}\)
    \item Ergio (erg): \(1 \, \text{erg} = 10^{-7} \, \text{J}\)
\end{itemize}

\section{Energía cinética}
\begin{equation}
    K = \frac{1}{2} m v^2
\end{equation}
\textbf{Teorema del trabajo y la energía:} El trabajo neto realizado sobre un cuerpo es igual al cambio en su energía cinética.
\begin{equation}
    W_{\text{neto}} = \Delta K = \frac{1}{2} m v_f^2 - \frac{1}{2} m v_i^2
\end{equation}

\section{Energía potencial}
\begin{align}
    \text{Gravitacional:} & U_g = m g h \\
    \text{Elástica:} & U_e = \frac{1}{2} k x^2
\end{align}

\section{Conservación de la energía mecánica}
Cuando solo actúan fuerzas conservativas:
\begin{equation}
    E = K + U = \text{constante}
\end{equation}

\section{Potencia}
\begin{equation}
    P = \frac{W}{\Delta t} = \vec{F} \cdot \vec{v}
\end{equation}
Unidad: Watt (\(1 \, \text{W} = 1 \, \text{J/s}\))

\resumebox{ejercicio}{Fuerza magnética del protón}{\footnotesize
    Un objeto de \(2 \, \text{kg}\) se eleva desde el suelo hasta una altura de \(5 \, \text{m}\). Calcula:
    \begin{enumerate}
        \item El trabajo realizado contra la gravedad.
        \item La energía potencial ganada.
        \item Si se suelta, ¿con qué velocidad llega al suelo?
    \end{enumerate}
}